\section{Introduction}
\label{sec:intro}

An AR device placing content on a table, or an indoor robot planning around furniture, can run a monocular metric depth foundation model~\cite{yang2024dav2,piccinelli2025unidepthv2,hu2024metric3dv2,bochkovskii2025depthpro} next to a cheap sparse depth sensor: a phone LiDAR, a time-of-flight module, a few-beam scanner or SLAM points.
Such a system needs, for each frame, an interval on every pixel's depth that holds for that frame, not on average over a dataset.
Conformal calibration inside the frame, on the sparse measurements, promises exactly this~\cite{mao2024spatial,gui2023cmc}, but only if the measurements are placed uniformly at random.
Real sensors return scan lines, clusters, depth edges, or only nearby surfaces.
When such structured measurements both correct the prediction and calibrate it (\cref{fig:teaser}), nearby measurements agree, the calibration sees small errors, and the intervals become \emph{narrower} while missing most of the frame: with clustered measurements, coverage drops to \NumLocalClusterCovLo--\NumLocalClusterCovHi{} at a $0.9$ target.
Only dense ground truth reveals this.

The sparse sensor is worth keeping.
Native confidences are not reliable~\cite{deng2026calibrationaudit}, and calibration across images gives only an average guarantee that hides images with far larger errors: an image's global scale error explains \NumScaleVarLo--\NumScaleVarHi\% of the image-to-image variation in interval coverage on NYU and KITTI.
Aligning to a few metric points~\cite{ranftl2022midas,wofk2023videpth} removes most of it: without anchors, split conformal prediction on a frozen model has several times the interval score of the anchored methods (\cref{sec:res:omni}).
We call this setting \emph{sparse-guided monocular depth}: a frozen foundation model supplies a dense prior, sparse metric measurements correct it, and calibration states what uncertainty remains.
It overlaps with depth completion~\cite{zuo2025omnidc,lin2025promptda}; our concern is the calibration step, and a completion network is one possible centre.

\paragraph{Idea.}
We give the sparse-depth budget two roles (\cref{fig:pipeline}).
\emph{Prediction anchors} follow the sensor layout, whatever it is, and build a local centre and scale around the frozen prediction.
\emph{Calibration probes}, a few measurements placed uniformly at random and never shown to the predictor, set one conformal threshold per image.
Because centre and scale are fixed before the probes are revealed, a rank argument gives a per-image coverage guarantee, in expectation over the probe draw, for any anchor layout and any centre that ignores the probes, including a depth-completion network, without calibration images.
Its spread across images follows an exact finite-population beta-binomial law, approximated by a Beta law, so the number of probes is a design choice that trades sensor budget for per-image reliability.
The guarantee certifies whatever the probes measure: for sensor readings, the readings.

The statistical ingredients are known: fitting and calibration sets inside one instance~\cite{gui2023cmc}, within-field exchangeability~\cite{mao2024spatial}, design-based conformal prediction~\cite{wieczorek2023design}, and the conditional coverage law of split conformal prediction~\cite{vovk2012conditional}.
Our contribution is their deployment-aware composition for sparse-guided dense depth, an analysis of the sensor budget, and a validation with an independent sensor:
\begin{itemize}[leftmargin=*,itemsep=1pt,topsep=2pt]
\item \textbf{Failure.} Structured sparse depth improves the prediction but, reused for calibration, voids the per-image guarantee; coverage falls short under scan-line, clustered and near-only layouts, with misses far from anchors (\cref{sec:res:failure}).
\item \textbf{A calibration layer.} Separating prediction anchors from calibration probes (\mDual, \cref{sec:method}) gives a per-image guarantee for any layout and any probe-independent centre; the probe budget sets the per-image spread, and a hull with distance-binned calibration adds a guarantee by distance to the nearest anchor (\mDualBoth). We evaluate three frozen models and five layouts (\cref{sec:res:dual,sec:res:budget}).
\item \textbf{Which centre, and what is certified.} On ARKitScenes~\cite{baruch2021arkitscenes} a depth-completion network~\cite{zuo2025omnidc} is the better centre under uniform anchors, the frozen model with our anchor-built centre under clustered ones (\cref{sec:res:omni}). With iPad LiDAR checked against laser scans, the layer certifies what the probes measure: laser probes certify the depth even with iPad anchors, while with iPad values for both, truth coverage falls far below the target under uniform anchors (\cref{sec:res:sensor}).
\end{itemize}
